\section{Numérique responsable}

La prise en compte du numérique responsable est principalement illustrée par le projet FCU / Ediacara. La réécriture des synchronisations évite de recopier les fichiers raster inchangés, limite les écritures sur les espaces partagés et centralise certaines livraisons réseau après un traitement local contrôlé. Lynx Immo apporte une preuve complémentaire de sobriété fonctionnelle par la définition d'un \terme{mvp} et le report de composants avancés tant que leur valeur et les données nécessaires n'étaient pas établies.

Ces éléments ne constituent pas une analyse environnementale complète. Ils montrent néanmoins l'intégration de critères d'efficience, de proportionnement et de limitation de la complexité dans les décisions techniques.
